% from ega4-18-fr.tex
\oldpage[IV-4]{188}

$\sh{O}_{Z,z}=A/\mathfrak R$ avec $\mathfrak J\subset\mathfrak R$. Par hypothèse, $\mathfrak J$ et $\mathfrak R$ sont des idéaux réguliers de $A$ et on a une suite exacte scindée
\[
  0\longrightarrow\mathfrak J/\mathfrak J\mathfrak R\longrightarrow\mathfrak R/\mathfrak R^2\longrightarrow(\mathfrak R/\mathfrak J)/(\mathfrak R/\mathfrak J)^2\longrightarrow0
\]
(cf. (16.2.7)). Cela entraîne qu'il existe des éléments $f_j$ ($1\leq j\leq r+s$) de $\mathfrak R$ dont les images $f'_j$ dans $\mathfrak R/\mathfrak R^2$ forment une base de cet $(A/\mathfrak R)$-module, et tels en outre que les $f_j$ d'indice $j\leq r$ appartiennent à $\mathfrak J$ et sont tels que leurs images dans $\mathfrak J/\mathfrak J\mathfrak R$ forment une base de cet $(A/\mathfrak R)$-module. Les $f_j$ pour $1\leq j\leq r+s$ engendrent donc $\mathfrak R$, et les $f_j$ pour $1\leq j\leq r$ engendrent $\mathfrak J$ puisque $A$ est noethérien (Bourbaki, \emph{Alg. comm.}, chap. II, \textsection~3, n\textsuperscript{o}~2, prop. 5); comme $\mathfrak R$ est un idéal régulier, il résulte de (16.9.5) que la suite $(f_j)_{1\leq j\leq r+s}$ est régulière. Par définition, les images de $f_{r+1},\ldots,f_{r+s}$ dans $\mathfrak R/\mathfrak J$ forment donc une suite régulière dans cet $(A/\mathfrak J)$-module, ce qui achève de prouver que $c)$ entraîne $a)$ dans (iv).

Enfin, compte tenu de (16.9.10), il est immédiat que (v), ainsi que l'équivalence de $a)$ et $b)$ dans (iv), résultent du

\begin{corollary}[19.1.6]
\label{IV.19.1.6}
Soient $A$ un anneau local noethérien, $\mathfrak J$, $\mathfrak R$ deux idéaux de $A$ contenus dans son idéal maximal $\mathfrak m$ et tels que $\mathfrak J\subset\mathfrak R$, $A'=A/\mathfrak J$, $\mathfrak R'=\mathfrak R/\mathfrak J$. Alors:
\begin{enumerate}
  \item[(i)] Si $\mathfrak J$ et $\mathfrak R'$ sont des idéaux réguliers, il en est de même de $\mathfrak R$.
  \item[(ii)] Si $\mathfrak R$ et $\mathfrak R'$ sont des idéaux réguliers, il en est de même de $\mathfrak J$.
  \item[(iii)] Supposons les anneaux $A$ et $A'=A/\mathfrak J$ réguliers. Pour que l'idéal $\mathfrak R$ soit régulier, il faut et il suffit que $\mathfrak R'$ le soit.
\end{enumerate}
\end{corollary}

L'assertion (i) est un cas particulier de (19.1.5, (iii)). Pour démontrer (ii), considérons l'homomorphisme surjectif canonique $u:\mathfrak R/\mathfrak R^2\to\mathfrak R'/\mathfrak R'^2=\mathfrak R/(\mathfrak R^2+\mathfrak J)$. Comme par hypothèse $\mathfrak R/\mathfrak R^2$ et $\mathfrak R'/\mathfrak R'^2$ sont des $(A/\mathfrak R)$-modules libres, dont nous désignerons les rangs respectifs par $p+q$ et $q$, le noyau de $u$ est un $(A/\mathfrak R)$-module \emph{projectif} (Bourbaki, \emph{Alg.}, chap. II, 3\textsuperscript{e} éd., \textsection~2, n\textsuperscript{o}~2, prop. 4), donc \emph{libre} de rang $p$ puisque $A/\mathfrak R$ est un anneau local noethérien ($0_{\mathrm{III}}$, 10.1.3); autrement dit, il existe une base de $\mathfrak R/\mathfrak R^2$ dont les $p$ premiers éléments forment une base de $(\mathfrak R^2+\mathfrak J)/\mathfrak R^2$. Cela signifie encore (en vertu du lemme de Nakayama) qu'il existe un système de générateurs $(f_i)_{1\leq i\leq p+q}$ de $\mathfrak R$, formant une suite régulière dans $A$, telle que la suite $(f_i)_{1\leq i\leq p}$ soit formée d'éléments de $\mathfrak J$; il s'agit de montrer que cette dernière suite engendre $\mathfrak J$. Pour cela, désignons par $\mathfrak J'\subset\mathfrak J$ l'idéal qu'elle engendre; en considérant l'anneau $A/\mathfrak J'$, on voit qu'on est ramené à prouver le lemme suivant:

\begin{lemma}[19.1.6.1]
\label{IV.19.1.6.1}
Soient $B$ un anneau local noethérien d'idéal maximal $\mathfrak n$, $(g_j)_{1\leq j\leq q}$ une suite régulière dans $B$ d'éléments de $\mathfrak n$, $\mathfrak b$ l'idéal qu'elle engendre, $\mathfrak c\subset\mathfrak b$ un idéal tel que les images canoniques des $g_j$ dans $B/\mathfrak c$ forment encore une suite régulière dans cet anneau. Alors on a $\mathfrak c=0$.
\end{lemma}

Le lemme étant trivial pour $q=0$, raisonnons par récurrence sur $q$. L'hypothèse de récurrence, appliquée à l'anneau $B/g_1B$ et aux images canoniques des $g_j$ ($2\leq j\leq q$)

